%% ===========================================================================
\section{Discussion}
\budget{1,300}

\subsection{Composition is the unit of analysis}
\budget{250}

\jun{Prose draft (my part; old bullet spec commented below). Covers the R8-vs-R2
reconciliation the spec flagged.}

The findings converge on a single methodological point: none of them is visible
in one annotation layer. The affective-privileging result needs sentiment,
frame, stance, and community together; the boundary-crossing result needs a
thirteen-dimensional author fingerprint; and the contestation result \emph{is},
in effect, the failure of the single-layer account, since the offensiveness layer
that a content score would read barely moves the outcome. The object that carries
the signal is the composition of layers with community context, not any layer on
its own.

This raises an apparent tension worth resolving directly, because a reader will
otherwise find it. Our contestation result says toxicity scarcely drives what a
community argues over, while the reply-tree analysis says frame-crossing
\emph{carries} more toxicity. Both hold, and they are not in conflict: toxicity
travels \emph{with} the act of crossing a frame boundary without \emph{driving}
whether a community contests a comment. The variable that predicts contestation
is the frame--community pairing; toxicity is a correlate of one particular move
within the discourse, not the quantity that decides its reception.

We state the scope of the claim carefully. This is a demonstration on one corpus
and one conflict that compositional structure recovers what a single layer hides;
it is not a proof that composition always dominates, and the mechanism behind each
composition remains to be tested on other conflicts and platforms.

% OLD SPEC (superseded by prose above):
% \begin{itemize}
% \item No headline finding is visible in one layer. R5 needs L2xL5xL4xcommunity;
%   R6 needs a 13-dim fingerprint; R8's result IS the single-layer failure.
% \item Reconcile: R8 says toxicity barely moves contestation; R2 says
%   frame-crossing carries toxicity. Both hold: toxicity travels WITH crossing
%   but does not DRIVE contestation.
% \item Generalization: one corpus, one conflict; not a proof composition always beats layer.
% \end{itemize}

\subsection{Humanitarian framing as the permeable seam}
\budget{250}

\jun{Prose draft (my part; old bullet spec commented below). Corrected the
audience-heterogeneity point: we tested it (§5.6) rather than leaving it
unresolved, and updated the comments-per-author figures to the verified values.}

Humanitarian framing is the seam along which these communities come apart. It is
the one frame that both diverges affectively across communities and, in the
contestation analysis, cracks consensus within them: it is simultaneously the
most contested category overall and the most likely to split an aligned community
that would close ranks around a security-and-military frame. Offered as a
mechanism and not a tested claim, the reading is that security-and-military
framing invites a rally-round-the-flag response an aligned community can answer in
unison, whereas humanitarian framing forces a moral judgment that co-partisans
answer differently --- so the disagreement surfaces inside the community rather
than only across the line.

The obvious competing explanation is that general-audience communities look more
contested simply because they draw a more transient crowd, whose one-off voters
mechanically split the vote --- an audience-heterogeneity effect rather than a
permeable echo chamber. Neutral communities do have the most transient audience
(2.69 comments per author, against 4.50 in india-aligned and 5.04 in
pakistan-aligned communities), so we tested the explanation rather than assert
against it. It does not survive: controlling for author activity leaves the
neutral effect essentially unchanged, and restricting to committed participants
\emph{widens} rather than shrinks the gap (\S5.6). What the data cannot rule out
is a subtler, voter-level version --- heterogeneity among the people who cast
votes, which the corpus does not observe --- and we mark that as the open
condition rather than claim to have closed it.

% OLD SPEC (superseded by prose above):
% \begin{itemize}
% \item The one frame that both diverges affectively (S5.1, S5.3) and cracks
%   consensus within communities (S5.6).
% \item Mechanism (not test): Sec/Mil = rally-round-the-flag; Humanitarian =
%   moral judgment co-partisans answer differently.
% \item Competing explanation "data cannot rule out" -> CORRECTED: tested in S5.6
%   and rejected (author transience); only voter-pool version remains open.
%   Numbers 2.73/4.64/5.21 -> 2.69/4.50/5.04.
% \end{itemize}

\subsection{``Bridge user'' is the wrong abstraction}
\budget{220}

\begin{itemize}
\item The clustering failure and the raiding signature together argue that
  crossing is a behaviour with a distribution, not a user type, and that its
  modal form is antagonistic.
\item Any system that treats cross-community participation as a proxy for
  constructive participation will systematically boost incursions.
\item Extends Garimella et al.~\cite{garimella2018echo} beyond network position;
  instantiates Kumar et al.~\cite{kumar2018conflict} at the individual level;
  sits against Xia et al.~\cite{xia2025integrated}, who find at best conditional
  depolarization from cross-party contact.
\end{itemize}

\subsection{Measurement as a first-class finding}
\budget{150}

\dnote{New in v2.}

\begin{itemize}
\item The specification curve's largest variance source is \textbf{the
  labeller}, not any sampling or modelling choice.
\item Perspective --- the classifier platforms actually deploy --- is the
  \textbf{weakest instrument measured} ($\kappa$ 0.228--0.249) against a human
  ceiling of 0.682 on the same construct, and was run English-only on a 3.1\%
  code-mixed corpus.
\item No automated labeller reached the human ceiling on any layer, and that
  ceiling is itself only 0.45--0.68.
\item Consequence for CSCW practice: LLM-annotated corpora need a
  \textbf{labeller-comparison arm}, not a single validation number --- and
  reliability reporting has to acknowledge that the human ruler is noisy too.
\end{itemize}

\subsection{Design implications}
\budget{280}

Four, each traced to a specific finding, each with its own failure mode named.

\begin{description}
\item[Composition-aware moderation triage.] Surface high-harm cells
  (Humanitarian $\rightarrow$ Security/Military transitions; Religious/Cultural
  $\times$ Offensive-Targeted in acute phases) rather than ranking by toxicity
  score. \emph{Failure mode:} composition cells are community-specific and will
  not transfer across conflicts without re-estimation.

\item[Narrative-transparency affordance.] Show per-community frame--affect
  composition on a thread. \emph{Failure mode:} making affective asymmetry
  legible may harden identity rather than soften it
  (cf.~\cite{bail2018exposure}).

\item[Crossing-aware visibility floors.] A minimum-visibility guarantee for
  out-group comments, since the penalty falls on inquisitive and balanced voices
  alike (question\_rate OR 2.15). \emph{Failure mode:} the same floor protects
  provocateurs, who are most of the crossers.

\item[Crisis-phase-adaptive policy.] Moderation thresholds as a function of
  \texttt{crisis\_phase} $\times$ \texttt{crisis\_type}, justified by \S5.1's
  finding that crisis type --- not era --- determines the discourse regime.
  \emph{Failure mode:} \S5.7 shows the corpus cannot causally identify phase
  effects, so this is a hypothesis for a platform-side experiment, not a
  validated recommendation.
\end{description}

\subsection{Few clusters and the limits of crisis-response causal inference}
\budget{150}

\jun{Prose draft (my part, C3; old bullet spec commented below).}

The methodological negative generalizes beyond this corpus. Any crisis-response
study whose treatment varies across roughly six events and ten communities sits
in the same few-clusters regime, where the cluster-robust variance estimator
over-rejects and the wild cluster bootstrap is only barely valid --- so an
uncorrected event-study on data of this shape will report effects that a proper
correction erases, as ours did on both designs (\S5.7, Figure~\ref{fig:crv1-wcb}).
The compounding problem is identification rather than power: when crises recur
inside each other's recovery windows, each event's pre-period is the aftermath of
the last, and parallel trends fail structurally. There are two ways out --- redesign
toward many clusters (thread-level opinion climates, author panels) or frame the
account as descriptive --- and this paper takes the second and says so plainly.

% OLD SPEC (superseded by prose above):
% \begin{itemize}
% \item Generalize C3: ~6 events x 10 communities is the regime where CRV1 lies
%   and WCB is barely valid.
% \item Two routes out: many-cluster redesigns or explicit descriptive framing.
%   This paper takes the second.
% \end{itemize}
